\chapter{Diseño Analítico de la Etapa en Surtidor Común}
\label{chap:diseno_teorico}

En este capítulo se expone el procedimiento sistemático de cálculo para dimensionar los componentes del circuito de
polarización en corriente continua, estableciendo criterios rigurosos de estabilidad térmica y optimización del rango
de excursión dinámica.

\section{Circuito Esquemático de Partida}

En la figura~\ref{crkt:surtidor_comun} se ilustra el esquema eléctrico completo del amplificador en surtidor común con
polarización por divisor de compuerta y resistencia de surtidor desacoplada.

\begin{figure}[!ht]
  \centering
  \begin{circuitikz}[scale=0.95, transform shape]
  % Fuente de alimentacion VDD
  \draw (3.5,6.5) node[vcc] {$V_{DD}$} coordinate (vdd_node);

  % Referencia y fuente de senal
  \draw (0,0) node[ground] {}
    to[sV, l=$v_s$] (0,3)
    to[C, l=$C_1$, -*] (3.5,3) coordinate (gate_node);

  % Divisor resistivo de compuerta (R1 a masa, R2 a VDD)
  \draw (vdd_node)
    to[R, l=$R_2$, -*] (gate_node)
    to[R, l=$R_1$, *-] (3.5,0) node[ground] {};

  % Transistor MOSFET Canal N
  \draw (gate_node) to[short] (4.5,3) node[nigfetd, anchor=G] (M1) {};
  \node[right=14pt, above=2pt] at (M1.right) {$M_1$};

  % Malla de Drenador con RD a VDD
  \draw (M1.D) to[short, -*] ++(0,0.8) coordinate (drain_node);
  \draw (drain_node) to[R, l=$R_D$] ++(0,2.0) coordinate (aux_d)
    to[short] (vdd_node -| aux_d) node[vcc] {$V_{DD}$};

  % Red de acoplo de salida desde el drenador hacia RL
  \draw (drain_node) to[C, l=$C_3$] ++(3.2,0) coordinate (out_node)
    to[short, -o] ++(0.8,0) node[above] {$v_L$};
  \draw (out_node) to[R, l=$R_L$, *-] ++(0,-3.8) node[ground] {};

  % Malla de Surtidor con RS y capacitor de desacoplo C2 en paralelo
  \draw (M1.S) to[short, -*] ++(0,-0.6) coordinate (source_node);
  \draw (source_node) to[R, l_=$R_S$] ++(0,-1.8) node[ground] {};
  \draw (source_node) to[short] ++(1.5,0)
    to[C, l=$C_2$] ++(0,-1.8) node[ground] {};
\end{circuitikz}

  \caption{Circuito esquemático completo del amplificador en configuración surtidor común con MOSFET 2N7000.}
  \label{crkt:surtidor_comun}
\end{figure}

El capacitor de desacoplo de entrada $C_1 = 100\uF$ aísla el nivel de continua de la compuerta respecto al generador
de funciones. El divisor resistivo integrado por $R_1$ y $R_2$ fija el potencial estático de compuerta ($V_G$). La
resistencia de surtidor $R_S$ provee realimentación negativa en corriente continua para estabilizar el punto de trabajo
frente a variaciones de temperatura y tolerancias del transistor. El capacitor $C_2 = 100\uF$ cortocircuita dinámicamente
a $R_S$ a la frecuencia de señal ($1\kHz$), derivando la corriente alterna hacia masa para preservar una elevada ganancia.
Finalmente, el capacitor $C_3 = 100\uF$ transfiere la componente senoidal de salida hacia la carga $R_L = 10\kohm$,
bloqueando el nivel continuo del drenador ($V_D$).

\section{Criterios de Diseño y Fijación del Punto de Operación}

La etapa se energiza mediante una fuente regulada estándar de corriente continua fijada en $V_{DD} = 15.0\V$. Se adoptan
los siguientes parámetros y especificaciones de partida:
\begin{itemize}
  \item Corriente de drenador en reposo: $I_{DQ} = 1.0\mA$.
  \item Parámetros medidos del transistor 2N7000: $K = 53.905\mA/\V^2$ y $V_{th} = 1.60104\V$.
  \item Impedancia de entrada del divisor: $R_G = R_1 \parallel R_2 = 100\kohm$.
  \item Criterio de estabilidad térmica de surtidor: se adopta una caída de tensión estática en surtidor equivalente al
    $10\%$ de la tensión de alimentación:
    \begin{equation}
      \frac{V_S}{V_{DD}} = 0.10 \implies V_S = 0.10 \cdot 15.0\V = 1.500\V
      \label{eq:criterio_VS}
    \end{equation}
\end{itemize}

\subsection{Cálculo de la Resistencia de Surtidor (\texorpdfstring{$R_S$}{RS})}
Aplicando la ley de Ohm en el terminal de surtidor:
\begin{equation}
  R_S = \frac{V_S}{I_{DQ}} = \frac{1.500\V}{1.0\mA} = 1.500\kohm = 1500\ohm
  \label{eq:RS_calculo}
\end{equation}

\subsection{Fijación de la Tensión de Drenador y Cálculo de \texorpdfstring{$R_D$}{RD}}
Con el fin de posibilitar la máxima excursión simétrica de tensión sobre el terminal de drenador sin aproximarse al corte
($V_D = V_{DD}$) ni a la tensión de surtidor ($V_S$), se sitúa el potencial de reposo en el punto medio del intervalo
disponible:
\begin{equation}
  V_D = \frac{V_{DD} + V_S}{2} = \frac{15.0\V + 1.500\V}{2} = \frac{16.500\V}{2} = 8.250\V
  \label{eq:VD_calculo}
\end{equation}

A partir de la malla de drenador, el resistor de carga de continua resulta:
\begin{equation}
  R_D = \frac{V_{DD} - V_D}{I_{DQ}} = \frac{15.0\V - 8.250\V}{1.0\mA} = \frac{6.750\V}{1.0\mA} = 6.750\kohm
  \label{eq:RD_calculo}
\end{equation}

La tensión continua drenador-surtidor de reposo queda establecida en:
\begin{equation}
  V_{DSQ} = V_D - V_S = 8.250\V - 1.500\V = 6.750\V
  \label{eq:VDSQ_calculo}
\end{equation}

\section{Diseño de la Red de Polarización de Compuerta}

A diferencia de un transistor bipolar, la corriente estática de compuerta en un MOSFET es prácticamente nula
($I_G \approx 0$) debido a la capa dieléctrica de $\mathrm{SiO}_2$. En consecuencia, el divisor formado por $R_1$ y
$R_2$ opera sin corriente de carga, comportándose como un divisor resistivo ideal.

\begin{figure}[!ht]
  \centering
  \begin{circuitikz}[scale=0.95, transform shape]
  % Fuente equivalente Thevenin VGG
  \draw (0,0) node[ground] {}
    to[battery2, l=$V_{GG}$] (0,2.5)
    to[R, l=$R_G$] (2.5,2.5) coordinate (gate_node);

  % Transistor MOSFET 2N7000
  \draw (gate_node) node[nigfetd, anchor=G] (M1) {};
  \node[right=14pt, above=2pt] at (M1.right) {$M_1$};

  % Drenador a VDD a traves de RD
  \draw (M1.D) to[R, l=$R_D$] ++(0,2.0) node[vcc] {$V_{DD}$};

  % Surtidor a masa a traves de RS
  \draw (M1.S) to[R, l=$R_S$] ++(0,-1.8) node[ground] {};
\end{circuitikz}

  \caption{Circuito equivalente de Thévenin para el análisis de polarización de compuerta.}
  \label{crkt:thevenin}
\end{figure}

Para lograr que circule la corriente de diseño $I_{DQ} = 1.0\mA$, se despeja la tensión compuerta-surtidor necesaria
a partir del modelo cuadrático de saturación:
\begin{equation}
  I_{DQ} = \frac{K}{2} \cdot (V_{GSQ} - V_{th})^2 \implies V_{GSQ} = V_{th} + \sqrt{\frac{2 \cdot I_{DQ}}{K}}
\end{equation}

Evaluando con los valores numéricos correspondientes:
\begin{equation}
  V_{GSQ} = 1.60104\V + \sqrt{\frac{2 \cdot 1.0\mA}{53.905\mA/\V^2}} = 1.60104\V + 0.19262\V \approx 1.79366\V \approx 1.794\V
  \label{eq:VGSQ_calculo}
\end{equation}

El potencial absoluto requerido sobre el nodo de compuerta respecto a masa resulta:
\begin{equation}
  V_G = V_S + V_{GSQ} = 1.500\V + 1.79366\V = 3.29366\V \approx 3.294\V
  \label{eq:VG_calculo}
\end{equation}

Dado que la compuerta no demanda corriente, la relación de transferencia de tensión del divisor viene dada por:
\begin{equation}
  \frac{V_G}{V_{DD}} = \frac{R_1}{R_1 + R_2} = \frac{3.29366\V}{15.0\V} \approx 0.219577
\end{equation}

Combinando esta relación con el valor adoptado de resistencia equivalente $R_G = R_1 \parallel R_2 = 100\kohm$:
\begin{equation}
  R_G = \frac{R_1 \cdot R_2}{R_1 + R_2} = R_2 \cdot \left(\frac{R_1}{R_1 + R_2}\right) = R_2 \cdot \left(\frac{V_G}{V_{DD}}\right)
\end{equation}
de donde se despejan las resistencias teóricas del divisor:
\begin{equation}
  R_2 = \frac{R_G}{V_G / V_{DD}} = \frac{100\kohm}{0.219577} \approx 455.420\kohm
  \label{eq:R2_teorico}
\end{equation}
\begin{equation}
  R_1 = \frac{R_G}{1 - (V_G / V_{DD})} = \frac{100\kohm}{1 - 0.219577} = \frac{100\kohm}{0.780423} \approx 128.136\kohm
  \label{eq:R1_teorico}
\end{equation}

\section{Verificación de Operación en Zona de Saturación}

Para que el modelo cuadrático mantenga estricta validez y la etapa se desempeñe como fuente de corriente controlada por
tensión (amplificador lineal de transconductancia), debe satisfacerse en todo instante:
\begin{equation}
  V_{DSQ} \ge V_{GSQ} - V_{th}
\end{equation}

Evaluando con las magnitudes de reposo calculadas:
\begin{equation}
  V_{DS,\mathrm{sat}} = V_{GSQ} - V_{th} = 1.79366\V - 1.60104\V = 0.19262\V \approx 0.193\V
\end{equation}
\begin{equation}
  V_{DSQ} = 6.750\V \gg 0.193\V
\end{equation}
comprobándose que el transistor se encuentra operando profundamente dentro de la zona de saturación, con un margen de
tensión de seguridad superior a $6.5\V$ respecto a la zona de triodo (óhmica).
